Se desarrolló una \textit{suite} de pruebas automatizadas que cubre las dos capas principales del sistema.

Las pruebas del \textit{backend} utilizan \texttt{django.test.TestCase} junto con el cliente \texttt{APIClient} de Django REST Framework. Los flujos asíncronos de \textit{streaming} SSE se cubren con \texttt{AsyncClient}, el cliente asíncrono nativo de Django, y la cobertura se mide con \texttt{coverage.py}. La \textit{suite} consta de \textbf{105 \textit{tests}} que verifican la autenticación, el ciclo de vida de las sesiones de chat, el \textit{streaming} del LLM, la moderación dinámica y la integración con Piston. Todos superan satisfactoriamente, alcanzando una cobertura del \textbf{99\,\%}.

Las pruebas del \textit{frontend} se ejecutan mediante \textbf{Vitest} en un entorno JSDOM, con \texttt{@vue/test-utils} para el montaje de componentes Vue y \texttt{@pinia/testing} para inyectar estados de \textit{store} controlados. En total se ejecutan \textbf{345 \textit{tests}}, todos satisfactorios, con una cobertura de sentencias del \textbf{91,6\,\%} global. Las vistas alcanzan el \textbf{93\,\%}, el \textit{sidebar} el \textbf{96\,\%} y los componentes de chat el \textbf{96,7\,\%}. El componente del editor de código se excluyó deliberadamente por su dependencia directa con Monaco~Editor, cuya validación requiere un entorno de navegador real.
